\documentclass[10pt,letterpaper,final]{report}
\usepackage[margin=1in]{geometry}
\usepackage[utf8]{inputenc}
\usepackage{amsmath}
\usepackage{amsfonts}
\usepackage{amssymb}
\usepackage{amsthm}
\renewcommand\qedsymbol{$\blacksquare$}
\usepackage{enumerate}
\usepackage{hyperref}
\usepackage{pdfpages}
\usepackage{graphics}
\usepackage{graphicx}
\usepackage{tikz}
\usepackage{tikz-qtree}
\usetikzlibrary{automata,arrows}
\usepackage{mdframed}
\usepackage[
  backend=biber,
  style=numeric-comp,
  sorting=none
]{biblatex}
\addbibresource{references.bib}

\usepackage{minted}
\usemintedstyle{colorful}
\usepackage{tabularx}


% Based on template from COSC-417 at Towson University, originally authored by Marius Zimand

\begin{document}
\begin{center}
  \fbox{
    \vbox{
      \begin{flushleft}
        Jonathan Llovet \\  % authors' names
        EN.605.202.81 \\  %class
        2026-06-16 \\  % date
        Module 4 - Homework 4 \\ % ID
      \end{flushleft}
      \center{\Large{\textbf{Data Structures - Queues and Lists}}}
      %\end{mdframed}
    } % end vbox
  } % end fbox
  \vline
\end{center}

Write pseudo-code for problems requiring code. Do not write Java, Python or C++. You are responsible for the appropriate level of detail.  For the questions asking for justification, please provide a detailed mathematically oriented discussion. A proof is not required.

\medskip

\noindent\textbf{Problem 1:}

Develop an ADT specification for a priority queue.  A priority queue is like a FIFO queue except that items are ordered by some priority setting instead of time. In fact, you may think of a FIFO queue as a priority queue in which the time stamp is used to define priority.

\medskip
\noindent\textbf{Answer:}
\medskip

\noindent\textbf{Priority Queue ADT}

\medskip

All references to \verb|q| in the following will be of type \verb|PriorityQueue|. That is, in the tables, \verb|q| can be read as the expression \verb|q: PriorityQueue|.

\noindent\textbf{Data}

\noindent\verb|q: PriorityQueue|

\begin{center}
  \begin{tabularx}{\textwidth}{|l|l|X|}
    \hline
    Field  & Type                  & Description                                                                                                                                                                                                   \\ \hline
    front  & \verb+a: Node | Null+ & Node at the front of the queue, i.e. the node with the highest priority                                                                                                                                       \\ \hline
    back   & \verb+b: Node | Null+ & Node at the back of the queue, i.e. the node with the lowest priority                                                                                                                                         \\ \hline
    length & \verb+x: int+         & Number of nodes in the queue. $x \in \mathbb{Z}\ \vert\ 0 \le x \le$ \verb|MAX_INT|, where \verb|MAX_INT| is the largest integer supported by the language or operating system or another external constraint \\ \hline
  \end{tabularx}
\end{center}

\noindent\verb|n: Node|

\begin{center}
  \begin{tabularx}{\textwidth}{|l|l|X|}
    \hline
    Field    & Type                   & Description                                                          \\ \hline
    data     & \verb|t: T|            & The data held by the node, here described as a generic type \verb|T| \\ \hline
    next     & \verb+m: Node* | Null+ & Pointer to the next node in \verb|q|                                 \\ \hline
    priority & \verb|p: int|          & Integer value used for determining position of the node in \verb|q|  \\ \hline
  \end{tabularx}
\end{center}

\pagebreak

\noindent\textbf{Operations}

In the table below, \emph{Purity} refers to whether there are side effects to a function, i.e. whether it modifies external state. A verb that is pure is idempotent and does not change external state. It receives a value as an input and usually produces a value as an output. A function that is not pure, in contrast, may produce a side effect, modifying state outside of itself.

\medskip
\begin{center}
  \begin{tabularx}{\textwidth}{|l|l|X|X|X|X|X|}
    \hline
    Operation       & Purity & Input    & Precondition              & Process                                                           & Postcondition & Output  \\ \hline
    \verb|is_empty| & Pure   & \verb|q| & \verb|q| must be non-null & Check length of \verb|q| and return whether it is currently empty & None          & Boolean \\ \hline
    \verb|enqueue|  &        & Input    & Precondition              & Process                                                           & Postcondition & Output  \\ \hline
    \verb|dequeue|  &        & Input    & Precondition              & Process                                                           & Postcondition & Output  \\ \hline
    \verb|contains| &        & Input    & Precondition              & Process                                                           & Postcondition & Output  \\ \hline
    \verb|iterate|  &        & Input    & Precondition              & Process                                                           & Postcondition & Output  \\ \hline
    \verb|remove|   &        & Input    & Precondition              & Process                                                           & Postcondition & Output  \\ \hline
    \verb|clear|    &        & Input    & Precondition              & Process                                                           & Postcondition & Output  \\ \hline
    \verb|promote|  &        & Input    & Precondition              & Process                                                           & Postcondition & Output  \\ \hline
    \verb|demote|   &        & Input    & Precondition              & Process                                                           & Postcondition & Output  \\ \hline
    \verb|shuffle|  &        & Input    & Precondition              & Process                                                           & Postcondition & Output  \\ \hline
  \end{tabularx}
\end{center}

\pagebreak

\noindent\textbf{Problem 2:}

Write an algorithm to reverse a singly linked list, so that the last element become the first and so on. Do NOT use Deletion - rearrange the pointers.

\medskip
\noindent\textbf{Answer:}

\pagebreak

\noindent\textbf{Problem 3:}

What is the average number of nodes accessed in search for a particular element in an unordered list? In an ordered list? In an unordered array? In an ordered array? Note that a list could be implemented as a linked structure or within an array.

\medskip
\noindent\textbf{Answer:}

\pagebreak

\noindent\textbf{Problem 4:}

Write a routine to interchange the mth and nth elements of a singly-linked list. You may assume that the ranks m and n are passed in as parameters. Allow for all the ways that m and n can occur. You must rearrange the pointers, not simply swap the contents.

\medskip
\noindent\textbf{Answer:}

\pagebreak

\section*{Source Code}
The full source code, including LaTeX, tooling, other artifacts, and git history is available on my Github repository for the class here:
\begin{center}
  \url{https://github.com/jllovet/jhu-en-605-202-su26-data-structures}
\end{center}

\printbibliography

\end{document}